\documentclass[10pt,a4paper]{article}
\usepackage[margin=1.55cm]{geometry}
\usepackage{xcolor}
\usepackage{amsmath,amssymb}
\usepackage{graphicx}
\usepackage{enumitem}
\usepackage{fontspec}
\usepackage[CJKmath=true]{xeCJK}
% CJK main font is resolved at COMPILE time on the actual machine, not baked in at
% generation time, so a .tex generated on one OS still renders on another (the older
% Python-side probe hardcoded one OS's font name into the file and broke when moved).
% Walk a cross-platform priority list and use the first font that exists; AutoFakeBold
% + AutoFakeSlant give bold/italic CJK even on faces that ship no bold/italic cut.
\newcommand{\setCJKto}[1]{\setCJKmainfont{#1}[AutoFakeBold=true,AutoFakeSlant=0.2]}
\IfFontExistsTF{Noto Sans CJK SC}{\setCJKto{Noto Sans CJK SC}}{%        % Linux (fonts-noto-cjk); cross-platform if installed
 \IfFontExistsTF{Source Han Sans SC}{\setCJKto{Source Han Sans SC}}{%   % Adobe 思源黑体; cross-platform
  \IfFontExistsTF{PingFang SC}{\setCJKto{PingFang SC}}{%               % macOS 10.11+
   \IfFontExistsTF{Hiragino Sans GB}{\setCJKto{Hiragino Sans GB}}{%    % older macOS
    \IfFontExistsTF{Microsoft YaHei}{\setCJKto{Microsoft YaHei}}{%     % Windows
     \IfFontExistsTF{WenQuanYi Micro Hei}{\setCJKto{WenQuanYi Micro Hei}}{% % older Linux
      \setCJKto{Songti SC}}}}}}}                                       % macOS bottom fallback
\definecolor{cblue}{HTML}{4C72B0}
\setlength{\parskip}{3pt}\setlength{\parindent}{0pt}
\setlist[enumerate]{topsep=2pt,itemsep=2pt,parsep=1pt,partopsep=0pt}
\setlist[itemize]{topsep=2pt,itemsep=2pt,parsep=1pt,partopsep=0pt}
% Let prose-embedded inline math break across lines and absorb small overflows, so simple
% formulas inserted mid-sentence stop running past the right margin.
\tolerance=1200
\emergencystretch=2.5em
\relpenalty=20
\binoppenalty=40
\hfuzz=1pt
\begin{document}

\section*{Boundary Score Matching for Functional-Geometry Counterfactuals}
\textbf{Method.} Boundary-Local Score Matching (BLSM)
\section*{Motivation}
A small change in clearance, pose, or contact geometry can change which actions are feasible even when the instruction, objects, and much of an action stay the same. Ordinary imitation gives a continuous-action vision-language-action (VLA) policy no direct signal for this boundary. Global feasibility supervision can reject an invalid rollout, but may spread the correction across the whole action chunk and lower the scores of nearby actions that remain physically valid. The key need is to assign the change to the first segment that actually violates the edited scene while retaining other valid modes.

Related systems provide geometry-aware validity, action manifolds, action-cluster selection, geometry-aware representations, or spatial-temporal action conditioning, but the available records do not combine a matched minimal edit, first-violation credit, complement anchoring, and mode-wise retention. Boundary-Local Score Matching (BLSM) makes the needed score change at a certified feasibility boundary distinct from score changes that should not occur elsewhere.

Recent VLA work and scripted RoboTwin 2.0 edits provide typed spatial grounding, persistent 3D context, geometry-guided validity, measurable action modes, paired counterfactuals, simulator-certified masks, and per-mode score accounting. A frozen continuous-action reference policy supplies the comparison needed for these score changes.

Closing the gap lets a VLA change the score for a particular contact or clearance segment after a geometry edit while retaining other feasible ways to finish the task. Reliability and diversity can then be measured together through invalid-action suppression, complement drift, alternative-positive coverage, and worst-mode success.
\section*{Method}
\subsection*{Background}
\emph{Construct accepted geometry-flip pairs, certify their first violation segment, and enumerate supported feasible modes.}
\begin{enumerate}[leftmargin=*,start=1]
  \item Create one structured record for every source rollout: task identifier, instruction, saved simulator state, the frozen policy's T-step continuous-action array and control rate, its original replay trace, and compiler version. First verify that the unchanged action array succeeds in the saved state. Search the declared task-preserving scene edits, replay that same array from the same state after each edit, and keep the edit with the smallest declared edit measure C, resolving ties deterministically. 【Author decision: state which scene parameters may change, their allowed domains and coupled constraints, how task identity is checked, how different edits are normalized in C, and how equal-C edits are resolved.】 Apply one versioned simulator rule F to label the complete original and edited replays. 【Author decision: define the success events, physical-invalidity events, numerical tolerances, duration requirements, and Boolean combination used by F for each task family.】 Save both scene states, the edit, C, action array, traces, tool versions, and acceptance status as the BoundaryFlip pair. Save every attempted edit in an eligibility ledger with its rule outcomes and a machine-readable inclusion or exclusion reason.
\noindent\emph{The compiler chooses the lowest-cost geometry edit that flips the identical action chunk from feasible to infeasible.}
\par\noindent{\small\ttfamily [eq~1, raw LaTeX, not typeset] \textbackslash\{\}Delta g\textasciicircum{}*=\textbackslash\{\}arg\textbackslash\{\}min\_\{\textbackslash\{\}Delta g\} C(\textbackslash\{\}Delta g)\textbackslash\{\}quad\textbackslash\{\}text\{s.t.\}\textbackslash\{\}quad F(x,a\textasciicircum{}-)=1,\textbackslash\{\};F(x\textbackslash\{\}oplus\textbackslash\{\}Delta g,a\textasciicircum{}-)=0}
\par\emph{Why:} Keeping the task and action identity fixed makes the geometry change the intended difference and shows whether eligible boundary pairs are broader than a selected corner of the data.
  \item Restore the pair's identical initial simulator state and random state, then replay the unchanged T-step action array in the original and edited scenes at the same control rate. After each action, log aligned state, contact, constraint, and task-event fields. Convert every edited-scene prefix into a binary violation value \(v_{t}.\) 【Author decision: define which task-specific events are monitored, their contact and state tolerances, how long an event must persist, when it clears, and how multiple events are combined into \(v_{t}.\)】 Find the earliest violating time \(\tau \) and continue through consecutive violating times until the first clear time or the end of the array. Set the first-violation mask \(m_{t}\) to one only on that run and set its complement on all other times. If no time violates, mark the pair mask-unresolved in the eligibility ledger instead of returning an empty mask. Store the two length-T bit arrays, \(\tau ,\) the run endpoint, supporting event values, and synchronized trace identifiers.
\noindent\emph{The mask is the first contiguous simulator-certified violation run beginning at the earliest violating time \(\tau .\)}
\par\noindent{\small\ttfamily [eq~2, raw LaTeX, not typeset] \textbackslash\{\}tau=\textbackslash\{\}min\textbackslash\{\}\{t:v\_t(x\textbackslash\{\}oplus\textbackslash\{\}Delta g\textasciicircum{}*,a\textasciicircum{}-\_\{1:t\})=1\textbackslash\{\}\},\textbackslash\{\}qquad m\_t=\textbackslash\{\}mathbf\{1\}\textbackslash\{\}!\textbackslash\{\}left[t\textbackslash\{\}in\textbackslash\{\}mathcal\{R\}(\textbackslash\{\}tau)\textbackslash\{\}right]}
\par\emph{Why:} The mask follows observed simulator structure and identifies the temporal region where the geometry edit first changes functional validity.
  \item Run the task compiler from the edited scene's saved initial state, using the same instruction, T-step horizon, control rate, and action representation as the invalid action array. Enumerate the compiler's declared solution domain, remove exactly duplicated action arrays, replay every candidate in the edited simulator, and retain only candidates accepted by the same complete-replay rule F. 【Author decision: specify the candidate representation, solver or generator, declared solution domain, duplicate rule, and exhaustion or coverage certificate used to call the enumeration complete.】 Group retained arrays with one deterministic definition of action mode. 【Author decision: define the object and equivalence rule that establish mode identity, including any trajectory, contact-sequence, object-side, path-class, or threshold criteria.】 Keep exactly one canonical scored array for every supported mode. 【Author decision: define how that representative is selected within a mode and the comparison measure, such as a compiler-designated canonical solution or a declared medoid.】 A mode is supported when at least one enumerated member passes F. Store the alternative-positive roster \(A^{+}\) in deterministic order;\allowbreak{} each row contains \(mode_{id},\) its T-by-D action array, compiler provenance, replay outcome, evidence for representative selection, the number of members in that mode, and total mode count K.
\par\emph{Why:} A mode-indexed roster makes positive coverage and worst-mode retention visible and prevents dominant-mode aggregation from hiding feasible alternatives.
\end{enumerate}
\subsection*{BLSM}
\emph{Apply boundary-local preference while anchoring the complement and every feasible alternative.}
\begin{enumerate}[leftmargin=*,start=4]
  \item Keep the pre-training reference checkpoint fixed in evaluation mode and start the trainable copy from that checkpoint. Under the edited observation, compute one scalar trainable score \(s_\theta \) and one scalar reference score \(s_{ref}\) at every action time for the invalid array and each mode representative. Use exactly the same observation preprocessing and conditioning context for both policies. 【Author decision: name the continuous-action policy form and define how its scalar score is calculated, including the conditioning context, reduction across action dimensions, normalization, and matching reference calculation.】 Before comparison, ensure that each alternative time index describes the same physical time as the first-violation mask \(m_{t}.\) 【Author decision: require every chunk to share T and the control rate, or define a deterministic alignment rule for unmatched times and phase-shifted contacts.】 Subtract the reference score from the matching trainable score to form the reference-adjusted change \(\Delta _\theta .\) For each mode, average within m the change of its feasible action relative to the invalid action to obtain margin \(d_\theta ^(k).\) Apply the specified smooth logistic preference to that margin and average all K modes equally, producing boundary loss \(L_{B}.\) Store \(L_{B},\) every mode identifier and margin, the number of selected times, and all per-time score records.
\noindent\emph{Every action score is expressed as a change from the frozen reference policy, preventing raw score scale from defining the preference.}
\par\noindent{\small\ttfamily [eq~3, raw LaTeX, not typeset] \textbackslash\{\}Delta\_\textbackslash\{\}theta(a\_t\textbackslash\{\}mid x')=s\_\textbackslash\{\}theta(a\_t\textbackslash\{\}mid x')-s\_\{\textbackslash\{\}mathrm\{ref\}\}(a\_t\textbackslash\{\}mid x')}
\noindent\emph{The boundary loss increases the masked reference-adjusted margin of each feasible alternative over the invalid action;\allowbreak{} \(M_{B}\) is the held-out mean margin.}
\par\noindent{\small\ttfamily [eq~4, raw LaTeX, not typeset] d\_\textbackslash\{\}theta\textasciicircum{}\{(k)\}=\textbackslash\{\}frac\{\textbackslash\{\}sum\_t m\_t[\textbackslash\{\}Delta\_\textbackslash\{\}theta(a\textasciicircum{}+\_\{k,t\}\textbackslash\{\}mid x')-\textbackslash\{\}Delta\_\textbackslash\{\}theta(a\textasciicircum{}-\_t\textbackslash\{\}mid x')]\}\{\textbackslash\{\}sum\_t m\_t\},\textbackslash\{\}quad \textbackslash\{\}mathcal\{L\}\_B=\textbackslash\{\}frac\{1\}\{K\}\textbackslash\{\}sum\_\{k=1\}\textasciicircum{}\{K\}\textbackslash\{\}log(1+e\textasciicircum{}\{-d\_\textbackslash\{\}theta\textasciicircum{}\{(k)\}\}),\textbackslash\{\}quad M\_B=\textbackslash\{\}mathbb\{E\}[d\_\textbackslash\{\}theta]}
\par\emph{Why:} Reference adjustment isolates post-training movement, while the mask directs feasibility credit to the counterfactual’s implicated segment rather than the whole chunk.
  \item For the invalid action array and all K mode representatives, take the times outside the first-violation mask, square each reference-adjusted score change \(\Delta _\theta ,\) average the values within each array, and then average the K+1 arrays equally. This produces complement loss \(L_{comp}.\) If the mask includes all T times, that average has no valid denominator. 【Author decision: state and record whether full-mask pairs are excluded, add no complement term, or use another explicitly defined treatment.】 Separately, square \(\Delta _\theta \) at every one of the T times for each feasible mode, average within each mode, and then average all K modes equally to produce alternative-preservation loss \(L_{alt}.\) Do not weight modes by how often the compiler generated them or by their member counts. Save both losses and the pair-level, mode-level, and time-level values needed to reproduce each reduction.
\noindent\emph{Complement anchoring limits off-boundary drift, while the equally weighted alternative loss anchors every feasible mode rather than only the dominant one.}
\par\noindent{\small\ttfamily [eq~5, raw LaTeX, not typeset] \textbackslash\{\}mathcal\{L\}\_\{\textbackslash\{\}mathrm\{comp\}\}=\textbackslash\{\}frac\{1\}\{K+1\}\textbackslash\{\}sum\_\{z\textbackslash\{\}in\textbackslash\{\}\{a\textasciicircum{}-,a\_1\textasciicircum{}+,\textbackslash\{\}ldots,a\_K\textasciicircum{}+\textbackslash\{\}\}\}\textbackslash\{\}frac\{\textbackslash\{\}sum\_t(1-m\_t)\textbackslash\{\}Delta\_\textbackslash\{\}theta(z\_t\textbackslash\{\}mid x')\textasciicircum{}2\}\{\textbackslash\{\}sum\_t(1-m\_t)\},\textbackslash\{\}quad \textbackslash\{\}mathcal\{L\}\_\{\textbackslash\{\}mathrm\{alt\}\}=\textbackslash\{\}frac\{1\}\{K\}\textbackslash\{\}sum\_\{k=1\}\textasciicircum{}\{K\}\textbackslash\{\}frac\{1\}\{T\}\textbackslash\{\}sum\_t\textbackslash\{\}Delta\_\textbackslash\{\}theta(a\textasciicircum{}+\_\{k,t\}\textbackslash\{\}mid x')\textasciicircum{}2}
\par\emph{Why:} The complement term prevents temporal credit leakage, and the alternative term makes suppressing minority feasible modes costly even when average validity improves.
\end{enumerate}
\subsection*{M3\_validation}
\emph{Select frozen loss weights and expose boundary, validity, drift, coverage, and worst-mode diagnostics.}
\begin{enumerate}[leftmargin=*,start=6]
  \item Build training batches from ordinary demonstrations and accepted BoundaryFlip records, and compute the base policy's usual imitation loss on its normal demonstration fields. 【Author decision: name the reference checkpoint and trainable policy class, action representation, conditioning inputs, initialization relation, native imitation loss, and batch reduction.】 Add boundary loss \(L_{B},\) complement loss \(L_{comp},\) and alternative-preservation loss \(L_{alt}\) using nonnegative weights \(\lambda _{B}, \lambda _{C},\) and \(\lambda _{A}.\) Use a validation split that keeps each original/edited pair together. Evaluate a predeclared Cartesian set of weight triples, keep only settings that meet the complement-drift and worst-mode non-degradation requirements, select the remaining setting with the greatest boundary separation, resolve ties deterministically, and then freeze the weights. 【Author decision: predeclare the candidate values for all three weights, the two non-inferiority margins, the validation statistic and uncertainty rule, and the tie breaker.】 Report held-out boundary margin \(M_{B}\) by averaging modes equally within each pair and then averaging pairs. Also report attempted and accepted eligibility counts with reasons, mask length as a count and fraction of T, roster size K, and held-out complement drift. Measure alternative-positive coverage and worst-mode retention by executing the trained policy in edited scenes. 【Author decision: define the rollout distribution, how executions are assigned to compiler modes, what makes a mode successfully covered, whether retention is absolute or relative to the reference policy, and how pairs and modes are aggregated.】 Save the complete training configuration, split manifest, selected weights, model checkpoint identifier, and pair-level diagnostic records. The output policy is a continuous-action Vision-Language-Action (VLA) policy trained with Boundary-Local Score Matching (BLSM).
\noindent\emph{BLSM adds boundary preference, complement stability, and mode preservation to the ordinary imitation objective;\allowbreak{} the three weights are selected once on validation pairs and then frozen.}
\par\noindent{\small\ttfamily [eq~6, raw LaTeX, not typeset] \textbackslash\{\}mathcal\{L\}=\textbackslash\{\}mathcal\{L\}\_\{\textbackslash\{\}mathrm\{imit\}\}+\textbackslash\{\}lambda\_B\textbackslash\{\}mathcal\{L\}\_B+\textbackslash\{\}lambda\_C\textbackslash\{\}mathcal\{L\}\_\{\textbackslash\{\}mathrm\{comp\}\}+\textbackslash\{\}lambda\_A\textbackslash\{\}mathcal\{L\}\_\{\textbackslash\{\}mathrm\{alt\}\}}
\par\emph{Why:} The selection rule prevents an arbitrary loss-weight trade-off from standing in for the mechanism, while the diagnostics connect score movement to reliability and retention.
\end{enumerate}
\end{document}
